\subsection{Continuity of certain representations}

PROPOSITION 1. — Let H be a locally compact group. Let \(\varrho_{1}\) (resp. \(\varrho_{2}\) ) be a continuous representation of G (resp. of H) in a locally convex separated K-vector space E \(_{1}\) (resp. E \(_{2}\) ). Let F denote the space \(\mathcal{L}(\mathrm{E}_{1};\mathrm{E}_{2})\) endowed with the topology of compact convergence. The representation \(\varrho\) of G×H into F defined by \(\varrho(g,h)u=\varrho_{2}(h)\circ u\circ\varrho_{1}(g^{-1})\) for \((g,h)\in\mathrm{G}\times\mathrm{H}\) is continuous.

Denote by \(\mathcal{L}_c(\mathrm{E}_1;\mathrm{E}_2)\) the space \(\mathcal{L}(\mathrm{E}_1;\mathrm{E}_2)\) equipped with the topology of compact convergence. Let \(\mathfrak{F}_1\) (resp. \(\mathfrak{F}_2,\mathfrak{F}_3\) ) be a filter in G converging to \(e\) (resp. a filter in \(\mathcal{L}_c(\mathrm{E}_1;\mathrm{E}_2)\) converging to 0, a filter in H converging to \(e\) ). Since G and H are locally compact, there exist elements \(\mathrm{C}\in \mathfrak{F}_1\) and \(\mathrm{D}\in \mathfrak{F}_3\) which are relatively compact. The set \(\varrho_1(\mathrm{C}^{-1})\) is equicontinuous in \(\mathcal{L}(\mathrm{E}_1)\) (cf. INT, VIII, p. 129, § 2, no. 1, remark 2, \(a'\) )). By prop. 9 of EVT, III, p. 33 and prop. 4 of EVT, III, p. 31, the map defined by \((u,v)\mapsto u\circ v\) of \(\mathcal{L}_c(\mathrm{E}_1;\mathrm{E}_2)\times \varrho_1(\mathrm{C}^{-1})\) in \(\mathcal{L}_c(\mathrm{E}_1;\mathrm{E}_2)\) is continuous. The filter base \(\mathfrak{F}_2\circ \varrho_1(\mathfrak{F}_1^{-1})\) therefore converges to 0 in \(\mathcal{L}_c(\mathrm{E}_1;\mathrm{E}_2)\) . The set \(\varrho_2(\mathrm{D})\) is equicontinuous in \(\mathcal{L}(\mathrm{E}_2)\) (cf. INT, VIII, p. 129, § 2, n° 1, rem. 2), and for every \(x\in \mathrm{E}_2\) , the set \(\varrho_2(\mathrm{D})x\subset \mathrm{E}_2\) is relatively compact. Consequently, \(\varrho_2(\mathrm{D})\) is relatively compact in \(\mathcal{L}(\mathrm{E}_2)\) endowed with the topology of compact convergence (TG, X, p. 18, cor. 1). The map defined by \((u,v)\mapsto u\circ v\) of \(\overline{\varrho_2(\mathrm{D})}\times \mathcal{L}_c(\mathrm{E}_1;\mathrm{E}_2)\) in \(\mathcal{L}_c(\mathrm{E}_1;\mathrm{E}_2)\) is continuous by prop. 9 of EVT, III, p. 33 and prop. 4 of EVT, III, p. 31. Hence the filter base \(\varrho (\mathfrak{F}_1\times \mathfrak{F}_3)(\mathfrak{F}_2)\) converges to 0 in \(\mathcal{L}_c(\mathrm{E}_1;\mathrm{E}_2)\) . This implies the assertion.

COROLLARY. --- Let \(\varrho\) be a continuous representation of G on a separated locally convex K-vector space E. The contragredient representation \(\breve{\varrho}\) is continuous when E' is endowed with the topology of compact convergence.

This follows from the proposition by taking \(H = \{e\}\) , \(\varrho_{1} = \varrho\) and \(\varrho_{2}\) the trivial representation on K.

Remark. --- The contragredient representation is not necessarily continuous when E' is endowed with the strong topology (cf. INT, VIII, p. 191, § 2, exercise 3, d)). One can show that it is continuous if the space E is semi-reflexive (loc. cit., c)).

